% Copia mecánica íntegra de bibliografia.tex:1--67.
% Residencia material del ensamblaje sucesor de 102 capítulos.
\renewcommand{\bibname}{Referencias de las correspondencias espirales genealógicas}
\begin{thebibliography}{99}

\bibitem{HaidaraRamosCierre}
O. Haidara Fall y R. Ramos Balsa,
\emph{El cierre holográfico del infinito: Holografía Modular Triádica y
Mecánica Dimensional}, edición integral, 2026.

\bibitem{HaidaraRamosMasa}
O. Haidara Fall y R. Ramos Balsa,
\emph{Teoría holográfica integral de la masa HMT--MD}, revisión 2, 2026.

\bibitem{HaidaraRamosContinuo}
O. Haidara Fall y R. Ramos Balsa,
\emph{La estructura discreta del continuo}, versión global 2, 2026.

\bibitem{ArnoldODE}
V. I. Arnold,
\emph{Ordinary Differential Equations},
Springer, tercera edición, 1992.

\bibitem{ArnoldMechanics}
V. I. Arnold,
\emph{Mathematical Methods of Classical Mechanics},
Springer, segunda edición, 1989.

\bibitem{Cartan}
É. Cartan,
\emph{Leçons sur la géométrie des espaces de Riemann},
Gauthier--Villars, 1928.

\bibitem{DoCarmo}
M. P. do Carmo,
\emph{Differential Geometry of Curves and Surfaces},
Prentice--Hall, 1976.

\bibitem{HilbertCohnVossen}
D. Hilbert y S. Cohn--Vossen,
\emph{Geometry and the Imagination},
Chelsea, 1952.

\bibitem{KatokHasselblatt}
A. Katok y B. Hasselblatt,
\emph{Introduction to the Modern Theory of Dynamical Systems},
Cambridge University Press, 1995.

\bibitem{MilnorDynamics}
J. Milnor,
\emph{Dynamics in One Complex Variable},
Princeton University Press, tercera edición, 2006.

\bibitem{Maxwell}
J. C. Maxwell,
\emph{A Treatise on Electricity and Magnetism},
Clarendon Press, 1873.

\bibitem{Weyl}
H. Weyl,
\emph{The Theory of Groups and Quantum Mechanics},
Dover, 1950.

\bibitem{LandauLifshitzFields}
L. D. Landau y E. M. Lifshitz,
\emph{The Classical Theory of Fields},
Butterworth--Heinemann, cuarta edición, 1975.

\end{thebibliography}
